\begin{table}[htbp]
\centering
\caption{E3 --- Sensitivity of the objective function to preference weight scaling on Semester~9 (10 replications each, Gurobi).}
\label{tab:e3_sensibilidad}
\scriptsize
\begin{tabular}{llrrrr}
\toprule
\textbf{Parameter} & \textbf{Scale factor} & \textbf{Obj.\ mean} & \textbf{Obj.\ std}
  & \textbf{$\Delta$Obj (\%)} & \textbf{Time (s) mean} \\
\midrule
  \multicolumn{2}{l}{\textit{Baseline} ($tn=md=ags=1.0$)} & 43.0 & 0.0 & 0 & 11.3 \\
  \midrule
  \multirow{3}{*}{$tn$ (shift--hour)} & 0.5 & 88.0 & 0.0 & +104.7 & 39.5 \\
   & 2.0 & 86.0 & 0.0 & +100.0 & 12.2 \\
   & 5.0 & 215.0 & 0.0 & +400.0 & 11.0 \\
  \midrule
  \multirow{3}{*}{$md$ (day preference)} & 0.5 & 88.0 & 0.0 & +104.7 & 37.5 \\
   & 2.0 & 86.0 & 0.0 & +100.0 & 10.9 \\
   & 5.0 & 215.0 & 0.0 & +400.0 & 10.4 \\
  \midrule
  \multirow{3}{*}{$ags$ (classroom preference)} & 0.5 & 88.0 & 0.0 & +104.7 & 39.6 \\
   & 2.0 & 86.0 & 0.0 & +100.0 & 9.8 \\
   & 5.0 & 215.0 & 0.0 & +400.0 & 11.2 \\
\botrule
\end{tabular}
\begin{tablenotes}
\small
\item $\Delta$Obj (\%) = percentage change in mean objective value relative to the baseline.
\item One parameter is scaled at a time; the other two are held at $1.0$.
\item Lower objective value indicates fewer idle gaps and better preference satisfaction.
\end{tablenotes}
\end{table}
